% Options for packages loaded elsewhere
\PassOptionsToPackage{unicode}{hyperref}
\PassOptionsToPackage{hyphens}{url}
\PassOptionsToPackage{dvipsnames,svgnames,x11names}{xcolor}
\documentclass[
  ngerman,
  10pt,
  a4paper,
]{article}
\usepackage{xcolor}
\usepackage[margin=22mm]{geometry}
\usepackage{amsmath,amssymb}
\setcounter{secnumdepth}{-\maxdimen} % remove section numbering
\usepackage{iftex}
\ifPDFTeX
  \usepackage[T1]{fontenc}
  \usepackage[utf8]{inputenc}
  \usepackage{textcomp} % provide euro and other symbols
\else % if luatex or xetex
  \usepackage{unicode-math} % this also loads fontspec
  \defaultfontfeatures{Scale=MatchLowercase}
  \defaultfontfeatures[\rmfamily]{Ligatures=TeX,Scale=1}
\fi
\usepackage{lmodern}
\ifPDFTeX\else
  % xetex/luatex font selection
  \setmainfont[Path=/Users/stefanhamann/.cache/codex-runtimes/codex-primary-runtime/dependencies/native/libreoffice-headless/libreoffice/LibreOfficeDev.app/Contents/Resources/fonts/truetype/,Extension=.ttf,BoldFont=DejaVuSans-Bold,ItalicFont=DejaVuSans-Oblique,BoldItalicFont=DejaVuSans-BoldOblique]{DejaVuSans}
  \setmonofont[Path=/Users/stefanhamann/.cache/codex-runtimes/codex-primary-runtime/dependencies/native/libreoffice-headless/libreoffice/LibreOfficeDev.app/Contents/Resources/fonts/truetype/,Extension=.ttf,BoldFont=DejaVuSansMono-Bold,ItalicFont=DejaVuSansMono-Oblique,BoldItalicFont=DejaVuSansMono-BoldOblique]{DejaVuSansMono}
\fi
% Use upquote if available, for straight quotes in verbatim environments
\IfFileExists{upquote.sty}{\usepackage{upquote}}{}
\IfFileExists{microtype.sty}{% use microtype if available
  \usepackage[]{microtype}
  \UseMicrotypeSet[protrusion]{basicmath} % disable protrusion for tt fonts
}{}
\makeatletter
\@ifundefined{KOMAClassName}{% if non-KOMA class
  \IfFileExists{parskip.sty}{%
    \usepackage{parskip}
  }{% else
    \setlength{\parindent}{0pt}
    \setlength{\parskip}{6pt plus 2pt minus 1pt}}
}{% if KOMA class
  \KOMAoptions{parskip=half}}
\makeatother
\usepackage{longtable,booktabs,array}
\newcounter{none} % for unnumbered tables
\usepackage{calc} % for calculating minipage widths
% Correct order of tables after \paragraph or \subparagraph
\usepackage{etoolbox}
\makeatletter
\patchcmd\longtable{\par}{\if@noskipsec\mbox{}\fi\par}{}{}
\makeatother
% Allow footnotes in longtable head/foot
\IfFileExists{footnotehyper.sty}{\usepackage{footnotehyper}}{\usepackage{footnote}}
\makesavenoteenv{longtable}
\ifLuaTeX
\usepackage[bidi=basic,shorthands=off]{babel}
\else
\usepackage[bidi=default,shorthands=off]{babel}
\fi
\ifPDFTeX
\else
\babelfont{rm}[Path=/Users/stefanhamann/.cache/codex-runtimes/codex-primary-runtime/dependencies/native/libreoffice-headless/libreoffice/LibreOfficeDev.app/Contents/Resources/fonts/truetype/,Extension=.ttf,BoldFont=DejaVuSans-Bold,ItalicFont=DejaVuSans-Oblique,BoldItalicFont=DejaVuSans-BoldOblique]{DejaVuSans}
\fi
\ifLuaTeX
  \usepackage{selnolig} % disable illegal ligatures
\fi
\setlength{\emergencystretch}{3em} % prevent overfull lines
\providecommand{\tightlist}{%
  \setlength{\itemsep}{0pt}\setlength{\parskip}{0pt}}
\usepackage{fvextra}
\RecustomVerbatimEnvironment{verbatim}{Verbatim}{breaklines=true,fontsize=\footnotesize}
\usepackage{xurl}
\usepackage{seqsplit}
\usepackage{adjustbox}
\usepackage{ragged2e}
\usepackage{titlesec}
\titleformat{\section}{\large\bfseries\raggedright}{}{0em}{}
\titleformat{\subsection}{\normalsize\bfseries\raggedright}{}{0em}{}
\DeclareRobustCommand{\codeword}[1]{\texttt{\seqsplit{#1}}}
\AtBeginEnvironment{longtable}{\small\RaggedRight}
\usepackage{newunicodechar}
\newunicodechar{𝒞}{\ensuremath{\mathcal C}}
\newunicodechar{𝔊}{\ensuremath{\mathfrak G}}
\newunicodechar{𝔤}{\ensuremath{\mathfrak g}}
\newunicodechar{𝓗}{\ensuremath{\mathcal H}}
\newunicodechar{ℋ}{\ensuremath{\mathcal H}}
\newunicodechar{𝔄}{\ensuremath{\mathfrak A}}
\newunicodechar{𝟙}{\ensuremath{\mathbf 1}}
\newunicodechar{𝕀}{\ensuremath{I}}
\newunicodechar{𝕋}{\ensuremath{\mathbb T}}
\newunicodechar{𝟘}{\ensuremath{\mathbf 0}}
\newunicodechar{⊞}{\ensuremath{\boxplus}}
\newunicodechar{∘}{\ensuremath{\circ}}
\newunicodechar{⟦}{\ensuremath{[\![}}
\newunicodechar{⟧}{\ensuremath{]\!]}}
\setlength{\emergencystretch}{4em}
\setlength{\parindent}{0pt}
\setlength{\parskip}{0.45em}
\AtBeginDocument{\hypersetup{colorlinks=true,linkcolor=black,urlcolor=blue}\pdfstringdefDisableCommands{\def\codeword#1{#1}}\RaggedRight}
\usepackage{fancyhdr}
\pagestyle{fancy}
\fancyhf{}
\fancyhead[L]{\small TFPT / Universalraum - v1.6.8}
\fancyfoot[R]{\thepage}
\setlength{\headheight}{15pt}
% The preserved full history now has three-digit page numbers.
\makeatletter
\renewcommand{\@pnumwidth}{2.6em}
\renewcommand{\@tocrmarg}{3.6em}
\makeatother
\usepackage{bookmark}
\IfFileExists{xurl.sty}{\usepackage{xurl}}{} % add URL line breaks if available
\urlstyle{same}
\hypersetup{
  pdflang={de-DE},
  colorlinks=true,
  linkcolor={Maroon},
  filecolor={Maroon},
  citecolor={Blue},
  urlcolor={Blue},
  pdfcreator={LaTeX via pandoc}}

\author{}
\date{}

\begin{document}

\section{TFPT / Universalraum: Kurzupdate
v1.6.8}\label{tfpt-universalraum-kurzupdate-v1.6.8}

\begin{enumerate}
\def\labelenumi{\arabic{enumi}.}
\setcounter{enumi}{14}
\tightlist
\item
  September 2026. Nur neue Befunde und Änderungen gegenüber v1.6.7.
  Vollständige Herleitungen und bewahrte Vorfassung stehen im
  Hauptdokument.
\end{enumerate}

\subsection{Die neue Forschungssynthese ist nützlich, aber nicht
fehlerfrei}\label{die-neue-forschungssynthese-ist-nuxfctzlich-aber-nicht-fehlerfrei}

Alle 74 archivierten Quellenhashes stimmen. Die 57 eigenen
Prüfbedingungen und die 20 übernommenen Kettenprüfungen wurden
unabhängig normal und optimiert wiederholt; die Ergebnisdateien stimmen
byteweise mit der Lieferung überein. Die große Grundzustandsprüfung und
die gesamte alte 822er-Suite wurden nicht erneut ausgeführt.

\textbf{Korrektur zu §8.4:} Vollständige Werte der ersten Gramzeile
reichen aus:

\[\adjustbox{max width=\linewidth}{$\displaystyle 
\Gamma(u,v)=\Gamma(1,u^\dagger v).
$}\]

Nur endlich viele ausgewählte Werte reichen im Allgemeinen nicht. Der
eigentliche GNS-Rekonstruktionssatz bleibt korrekt. Ein markierter
komplexer Wortkern setzt außerdem mehr Information voraus als
gewöhnliche Quantenkanäle ohne kohärente Zweigkontrolle.

Die Ein-Cartan-Referenz und die bewegte Projektorgeometrie des Papers
sind korrekte bedingte Konstruktionen. Referenz, Kopplung und
Projektorfamilie sind keine bereits abgeleiteten nativen Ressourcen. In
der ursprünglichen reellen hellen Zweikomponentenfamilie verschwindet
die glatte lokale Berry-Zweiform.

\subsection{Drei parallele
Fortschritte}\label{drei-parallele-fortschritte}

\textbf{1. Ein kleiner ursprünglicher Mechanismus.} Zwei volle innere
Singuletts bei N=4 schließen exakt unter dem gesamten nativen H:

\[\adjustbox{max width=\linewidth}{$\displaystyle 
H_4=\begin{pmatrix}2\Delta&4g\\4g&\Delta\end{pmatrix}.
$}\]

Die Fierz-Identität wurde an allen 3840 signierten Beiträgen geprüft. Am
Prüfpunkt g/Δ=1/20 ist die maximale Umwandlung 4/29. Vorbereitung aus
dem Vakuum bleibt offen. Kein Wort des dokumentierten passiven
Operationsalphabets ändert N. Ein zusätzlich gewährter Bosonpaargriff Q
erzeugt den zweiten Erweiterungskanal algebraisch durch
\(R_++R_-=-[N_b,[Q,X]]\). Eine neue zweistufige Referenz bilanziert ihn
unter voller G-Symmetrie; sie ist kein kostenloser isolierter Antrieb.

\textbf{2. Eine vollständige bedingte Schattenkarte.} Sichtbare
Operatorränge im inneren 60-Raum:

{\def\LTcaptype{none} % do not increment counter
\begin{longtable}[]{@{}lr@{}}
\toprule\noalign{}
Messfamilie & Rang \\
\midrule\noalign{}
\endhead
\bottomrule\noalign{}
\endlastfoot
Einzelbesetzungen & 24 \\
Alle Bilineare & 736 \\
Bilineare plus ungedrehte Paare & 772 \\
G-gedrehte Paarmessungen & 3600 \\
\end{longtable}
}

Eine explizite Familie von 3599 unabhängigen nichtkonstanten
Erwartungswerten rekonstruiert jeden normierten inneren Zustand. Acht
unbenutzte Zustände und 96 zurückgehaltene Vorhersagen bestätigen die
Rekonstruktion numerisch. Ausführbare Gruppenkontrollen und gemeinsame
Paarmessungen bleiben Voraussetzungen.

\textbf{Neu verbindend:} Beide Zweigbilder zu drei Zeiten bestimmen auch
die relative Paar-/Bosonphase. Im hellen Raum gilt exakt

\[\adjustbox{max width=\linewidth}{$\displaystyle 
\dot S=0,\quad\dot X=\Delta Y,\quad
\dot Y=-\Delta X-2aZ,\quad\dot Z=2aY,
\quad a=\sqrt8g.
$}\]

Für die Zeiten mit \(\omega t=0,\pi/2,\pi\) ist die
Rekonstruktionsdeterminante \(8a^2\Delta/\omega^3\ne0\). Ein vierter
Zeitpunkt wird ohne Anpassung exakt vorhergesagt. Nur Paardaten bleiben
selbst über alle Zeiten unvollständig; die zweite
Zweigwahrscheinlichkeit darf nicht wegnormiert werden.

\textbf{3. Ein echter kausaler Modeneingriff.} Im ursprünglichen N=2-H
verändert eine gezielte lokale Phase an Mode 4 die spätere unbedingte
Besetzung von Mode 5. Kein zusätzlicher Hoppingterm, keine
Nachselektion. Beide Arme enden exakt mit zwei Fermionen und null
Bosonen. Mit

\[\adjustbox{max width=\linewidth}{$\displaystyle 
\Omega^2=\Delta^2/4+8g^2,\quad T=\pi/\Omega,
\quad x=4\cos^2\big(\pi\Delta/(4\Omega)\big)
$}\]

ist die Differenz nach Vorentwicklung, Impuls und Nachentwicklung

\[\adjustbox{max width=\linewidth}{$\displaystyle 
\delta_{4\to5}=x(25x-64)/1024.
$}\]

Sie ist im angegebenen Parameterintervall strikt negativ, am Prüfpunkt
etwa -0.00087318. Präparation, gezielte Phase und Auslesung sind
weiterhin gewährte Instrumente, keine hergeleitete Compilerfähigkeit.

\textbf{Zusammenführung auf demselben Träger:} Der kausale Versuch
schließt bereits auf drei Zuständen: Anfangspaar, kohärente Summe der
sieben anderen Paare und Boson. Auf genau diesem Raum bestimmen
terminale Besetzungsdaten samt zeitlichen Ableitungen den ganzen
Zustand. Die symbolische Observablendeterminante ist
\(8\Delta^3g^{10}/343\ne0\). Damit sind Rekonstruktion und Eingriff in
einer kleinen nativen Ausführung verbunden. Die letzte
Einzelmodenmessung kann jedoch aus diesem Raum herausführen;
Wiederverwendung nach dieser Messung ist nicht Teil des Satzes.

\subsection{Noch nachgereicht: zwei zusätzliche wichtige
Befunde}\label{noch-nachgereicht-zwei-zusuxe4tzliche-wichtige-befunde}

Die fünfte Bosonstufennorm \(\nu_5=1866738327552000\) ist frisch mit
exakter ganzzahliger Spurkontraktion reproduziert. Ein unabhängig
zertifizierter Siebener-Variationsraum liefert
\(E_0<-1.13847609\Delta\). Mit dem beibehaltenen N=63-Satz folgt
\(\epsilon>0.01657709\Delta\), statt zuvor 0.00773911Δ. Das ist keine
geschlossene Siebener-Lanczoskette und kein Polzentralwert.

Die einfache Quelle vermischt Zahlen aus einer Fünfer- und einer
Siebener-Näherung. Beispielsweise gehören 0.973684/0.026316 nicht zu dem
größeren Zustand mit Energie um -1.138476Δ. Dessen gesamte
Entnahme-/Additionsnormen liegen näherungsweise bei 0.9681955/0.0318045;
auch das sind keine exakten Polgewichte auf Ω.

Der behauptete Spin-1/2-Projektor ist im gleichhändigen
Tensorproduktraum korrekt. Bei einem wörtlich adjungierten Weylfeld
entsteht dagegen ein gemischthändiger Lorentzraum: Der
Rotationsprojektor kommutiert nicht mit Boosts, sein entsprechender
Kommutator hat Rang vier. Ein korrekt typisierter, möglicherweise
unabhängiger oder ladungskonjugierter Feldadapter bleibt deshalb offen.
Das Symbol f† allein legt seine Händigkeit nicht fest.

Die dritte Notiz bestätigt: Wir müssen die globale Zusammensetzung aus
der Quelle herleiten. Sie liefert keinen neuen Mechanismus und ersetzt
nicht die fehlende Ableitung räumlicher Vorkommen. Die interne
Zustandspräparation ist von einer globalen Zustandsregel ohne äußeren
Experimentator zu unterscheiden.

\subsection{Was sich nicht geändert
hat}\label{was-sich-nicht-geuxe4ndert-hat}

Die neuen N=2-/N=4-Zeugen ersetzen nicht den N=64-Grundzustand oder
seinen geladenen Pol. Ein endlicher kausaler Modeneffekt ist keine
hergeleitete räumliche Ausbreitung. Alle vollständigen T1-T8-Tore
bleiben offen. Jetzt ist aber ein konkreter kleiner Instrumentenvertrag
vorhanden, an dem sich der nächste Herkunftsnachweis entscheiden lässt.

\end{document}
